% ============================================================
%DIF LATEXDIFF DIFFERENCE FILE
%DIF DEL original.tex                    Fri Oct 16 16:11:31 2026
%DIF ADD primera_entrega_corregida.tex   Fri Oct 16 16:11:31 2026
%  Desigualdad del ingreso laboral en Colombia, 2018-2026
%  Formato APA 7 (estudiante) -- compila con pdfLaTeX
% ============================================================
\documentclass[12pt,letterpaper]{article}

% ---------- Idioma y fuentes ----------
\usepackage[T1]{fontenc}
\usepackage[utf8]{inputenc}
\usepackage[spanish,es-nodecimaldot,es-nolists,es-noshorthands]{babel}
\usepackage{lmodern}
\usepackage{amsmath,amssymb}
\usepackage{icomma}            % permite escribir 0,05 en modo matemático
%DIF 14a14
\usepackage{graphicx}
\usepackage{float}          % logo de la portada (ur.png en la misma carpeta) %DIF > 
%DIF -------

% ---------- Página APA ----------
\usepackage[margin=1in]{geometry}
\usepackage{setspace}
\doublespacing
\setlength{\parindent}{0.5in}
\setlength{\parskip}{0pt}

\usepackage{fancyhdr}
\pagestyle{plain}
\fancyhf{}
\fancyhead[R]{\thepage}
\renewcommand{\headrulewidth}{0pt}
\setlength{\headheight}{15pt}

% ---------- Títulos APA (sin numeración) ----------
\usepackage{titlesec}
\setcounter{secnumdepth}{0}
% Nivel 1: centrado, negrita
\titleformat{\section}[block]{\normalfont\bfseries\centering}{}{0pt}{}
\titlespacing*{\section}{0pt}{0pt}{0pt}
% Nivel 2: alineado a la izquierda, negrita
\titleformat{\subsection}[block]{\normalfont\bfseries}{}{0pt}{}
\titlespacing*{\subsection}{0pt}{0pt}{0pt}

% ---------- Listas, tablas y enlaces ----------
\usepackage{enumitem}
\setlist{nosep}
\usepackage{array,booktabs}
\usepackage{xurl}
\usepackage[hidelinks]{hyperref}
\urlstyle{same}

%DIF 47-54c48-57
%DIF < % Lista de referencias con sangría francesa
%DIF < \newenvironment{referencias}{%
%DIF <   \begin{list}{}{%
%DIF <     \setlength{\leftmargin}{0.5in}%
%DIF <     \setlength{\itemindent}{-0.5in}%
%DIF <     \setlength{\itemsep}{0pt}%
%DIF <     \setlength{\parsep}{0pt}}}%
%DIF <   {\end{list}}
%DIF -------
% ---------- Citación APA 7 (biblatex + biber) ---------- %DIF > 
% Compilar: pdflatex -> biber -> pdflatex -> pdflatex %DIF > 
% (Overleaf lo hace solo; el archivo references.bib debe estar en la misma carpeta) %DIF > 
\usepackage{csquotes} %DIF > 
\usepackage[backend=biber,style=apa]{biblatex} %DIF > 
\addbibresource{references.bib} %DIF > 
% APA en español: "y" entre autores en lugar de "&" %DIF > 
\DeclareDelimFormat[bib,biblist]{finalnamedelim}{% %DIF > 
  \ifnum\value{liststop}>2 \finalandcomma\fi\addspace y\space} %DIF > 
\DeclareDelimFormat[parencite]{finalnamedelim}{\addspace y\space} %DIF > 
%DIF -------


\newcommand{\titulo}{¿Cómo estimar el índice de Theil con la GEIH? Comparación de cuatro estimadores aplicado a la desigualdad del ingreso laboral en Colombia, 2018--2026}

% ============================================================
%DIF PREAMBLE EXTENSION ADDED BY LATEXDIFF
%DIF PREAMBLE: azul = texto nuevo; rojo tachado = texto anterior
\RequirePackage[normalem]{ulem} %DIF PREAMBLE
\RequirePackage{color}\definecolor{RED}{rgb}{0.8,0,0}\definecolor{BLUE}{rgb}{0,0,0.85} %DIF PREAMBLE
\providecommand{\DIFaddtex}[1]{{\protect\color{BLUE}#1}} %DIF PREAMBLE
\providecommand{\DIFdeltex}[1]{{\protect\color{RED}\sout{#1}}} %DIF PREAMBLE
\providecommand{\DIFaddtex}[1]{\texorpdfstring{\DIFaddtex{#1}}{#1}} %DIF PREAMBLE
\providecommand{\DIFdeltex}[1]{\texorpdfstring{\DIFdeltex{#1}}{}} %DIF PREAMBLE
\providecommand{\DIFaddbegin}{}\providecommand{\DIFaddend}{} %DIF PREAMBLE
\providecommand{\DIFdelbegin}{}\providecommand{\DIFdelend}{} %DIF PREAMBLE
\providecommand{\DIFmodbegin}{}\providecommand{\DIFmodend}{} %DIF PREAMBLE
\providecommand{\DIFaddFL}[1]{\DIFadd{#1}}\providecommand{\DIFdelFL}[1]{\DIFdel{#1}} %DIF PREAMBLE
\providecommand{\DIFaddbeginFL}{}\providecommand{\DIFaddendFL}{} %DIF PREAMBLE
\providecommand{\DIFdelbeginFL}{}\providecommand{\DIFdelendFL}{} %DIF PREAMBLE
%DIF END PREAMBLE %DIF PREAMBLE
%DIF HYPERREF PREAMBLE %DIF PREAMBLE
\providecommand{\DIFadd}[1]{\texorpdfstring{\DIFaddtex{#1}}{#1}} %DIF PREAMBLE
\providecommand{\DIFdel}[1]{\texorpdfstring{\DIFdeltex{#1}}{}} %DIF PREAMBLE
%DIF COLORLISTINGS PREAMBLE %DIF PREAMBLE
\RequirePackage{listings} %DIF PREAMBLE
\RequirePackage{color} %DIF PREAMBLE
\lstdefinelanguage{DIFcode}{ %DIF PREAMBLE
%DIF DIFCODE_UNDERLINE %DIF PREAMBLE
  moredelim=[il][\color{red}\sout]{\%DIF\ <\ }, %DIF PREAMBLE
  moredelim=[il][\color{blue}\uwave]{\%DIF\ >\ } %DIF PREAMBLE
} %DIF PREAMBLE
\lstdefinestyle{DIFverbatimstyle}{ %DIF PREAMBLE
	language=DIFcode, %DIF PREAMBLE
	basicstyle=\ttfamily, %DIF PREAMBLE
	columns=fullflexible, %DIF PREAMBLE
	keepspaces=true %DIF PREAMBLE
} %DIF PREAMBLE
\lstnewenvironment{DIFverbatim}{\lstset{style=DIFverbatimstyle}}{} %DIF PREAMBLE
\lstnewenvironment{DIFverbatim*}{\lstset{style=DIFverbatimstyle,showspaces=true}}{} %DIF PREAMBLE
%DIF END PREAMBLE EXTENSION ADDED BY LATEXDIFF

\begin{document}

% ---------------- Portada ----------------
\DIFdelbegin %DIFDELCMD < \begin{titlepage}
%DIFDELCMD < \thispagestyle{empty}
%DIFDELCMD <   \centering
%DIFDELCMD <   \vspace*{2in}
%DIFDELCMD <   {\bfseries \titulo\par}
%DIFDELCMD <   \vspace{\baselineskip}
%DIFDELCMD <   Gabriel Salcedo, Jerónimo Marín, Isaac Suarez y Santiago Velandia\par
%DIFDELCMD <   Universidad del Rosario\par
%DIFDELCMD <   Estadistica\par
%DIFDELCMD <   Octubre de 2026\par
%DIFDELCMD < \end{titlepage}
%DIFDELCMD < %%%
\DIFdelend \DIFaddbegin \begin{titlepage}
  \thispagestyle{empty}
  \centering
  \singlespacing
  \vspace*{0.2in}
  \includegraphics[width=2.4in]{ur.png}\par
  \vspace{0.9in}
  \rule{0.75\textwidth}{0.4pt}\par
  \vspace{0.3in}
  {\large\bfseries \titulo\par}
  \vspace{0.3in}
  \rule{0.75\textwidth}{0.4pt}\par
  \vspace{0.9in}
  Gabriel Salcedo\par
  Jerónimo Marín\par
  Isaac Suarez\par
  Santiago Velandia\par
  \vfill
  Universidad del Rosario\par
  \vspace{0.15in}
  Estadística\par
  \vspace{0.15in}
  Octubre de 2026\par
  \vspace*{0.2in}
\end{titlepage}
\DIFaddend \setcounter{page}{2}

% ---------------- Leyenda de cambios ----------------
\begin{center}
\fbox{\begin{minipage}{0.85\textwidth}\singlespacing\small
\textbf{Versión con control de cambios.} Compara la primera entrega con la versión corregida.\\[3pt]
{\color{BLUE}Texto en azul}: texto nuevo o corregido.\quad
{\color{RED}\sout{Texto rojo tachado}}: texto de la versión anterior.\\[3pt]
Las tablas que cambiaron aparecen completas dos veces: la anterior en rojo y la corregida en azul. Las citas ahora salen de \texttt{references.bib}; por eso las referencias aparecen dos veces: la lista anterior tachada y la nueva, generada automáticamente.
\end{minipage}}
\end{center}


% ---------------- Título en la primera página del texto ----------------

\section{Introducción}

\subsection{Justificación}

%DIF >  REVISAR: el boletín de pobreza monetaria 2024 (clave dane2025pobreza) no parece reportar el Gini de 2000.
%DIF >  Confirmar la cifra y su fuente, o volver a la comparación con 2023 (0,553).
Colombia es uno de los países más desiguales de América Latina: el Gini del ingreso de los hogares fue 0,551 en 2024, casi igual al 0,545 de 2000 \DIFdelbegin \DIFdel{(DANE, 2025)}\DIFdelend \DIFaddbegin \parencite{dane2025pobreza}\DIFaddend . Estas cifras se calculan a partir de encuestas, y por eso no son la desigualdad ``real'' sino una estimación de ella. \DIFdelbegin \DIFdel{Por eso, detrás }\DIFdelend \DIFaddbegin \DIFadd{Detrás }\DIFaddend de cada dato publicado hay decisiones \DIFdelbegin \DIFdel{estadisticas }\DIFdelend \DIFaddbegin \DIFadd{estadísticas }\DIFaddend sobre qué \DIFdelbegin \DIFdel{estimadores se usan }\DIFdelend \DIFaddbegin \DIFadd{estimador se usa }\DIFaddend y bajo qué supuestos.

El índice de Theil es una medida usual para este problema, pues \DIFdelbegin \DIFdel{, }\DIFdelend permite calcular la desigualdad \DIFdelbegin \DIFdel{de ingreso yademas es capaz de  descomponer el dato de }\DIFdelend \DIFaddbegin \DIFadd{del ingreso y, además, descomponerla exactamente en }\DIFaddend desigualdad \emph{entre} grupos y \emph{dentro} de ellos \DIFdelbegin \DIFdel{(Shorrocks, 1980)}\DIFdelend \DIFaddbegin \parencite{shorrocks1980class}\DIFaddend . Pero hay varias formas de estimarlo con los mismos datos. Se puede calcular directamente sobre la muestra, con o sin los factores de expansión de la encuesta, o se puede suponer que el ingreso sigue una distribución lognormal, \DIFdelbegin \DIFdel{lo que es usual en la distribuciones de ingreo}\DIFdelend \DIFaddbegin \DIFadd{un supuesto usual en las distribuciones de ingreso}\DIFaddend . Cada opción necesita supuestos distintos, y elegir mal puede producir sesgo.

En este proyecto queremos comparar, mediante un experimento de simulación, cuatro estimadores del índice de Theil en términos de sesgo, varianza, error cuadrático medio (ECM), eficiencia relativa y consistencia\DIFdelbegin \DIFdel{. Y }\DIFdelend \DIFaddbegin \DIFadd{, y }\DIFaddend después usar el estimador que mejor se comporte para medir la desigualdad del ingreso laboral en Colombia entre 2018 y 2026.

\subsection{Estado del arte}

\DIFdelbegin \DIFdel{Theil (1967) }\DIFdelend \DIFaddbegin \DIFadd{\textcite{theil1967economics} }\DIFaddend propuso un índice que vale cero cuando todas las personas ganan lo mismo y crece a medida que el ingreso se concentra. Este \DIFdelbegin \DIFdel{indice posee }\DIFdelend \DIFaddbegin \DIFadd{índice tiene }\DIFaddend varias propiedades que nos serán \DIFdelbegin \DIFdel{utiles mas adelante, como: }\DIFdelend \DIFaddbegin \DIFadd{útiles más adelante; la principal es }\DIFaddend que es aditivamente descomponible, y a partir de eso se puede medir \DIFdelbegin \DIFdel{cuanto explica cada división (aditiva) el valor global del Theil }\DIFdelend \DIFaddbegin \DIFadd{cuánto del Theil total explica cada partición de la población en grupos}\DIFaddend .

Por otro lado, para cuantificar la incertidumbre de cualquier estimador, \DIFdelbegin \DIFdel{Efron y Tibshirani (1993) }\DIFdelend \DIFaddbegin \DIFadd{\textcite{efron1993introduction} }\DIFaddend desarrollaron el bootstrap, \DIFdelbegin \DIFdel{él cual es un }\DIFdelend un método que aproxima la distribución muestral de un estimador remuestreando los datos, y \DIFdelbegin \DIFdel{Mills y Zandvakili (1997) }\DIFdelend \DIFaddbegin \DIFadd{\textcite{mills1997statistical} }\DIFaddend lo aplicaron a índices de desigualdad y a sus descomposiciones.

Para Colombia, \DIFdelbegin \DIFdel{Sánchez-Torres (2017) }\DIFdelend \DIFaddbegin \DIFadd{\textcite{sanchez2017desigualdad} }\DIFaddend encontró que cerca del 88~\% de la desigualdad se da dentro de los departamentos, y \DIFdelbegin \DIFdel{Ariza y Montes-Rojas (2017) }\DIFdelend \DIFaddbegin \DIFadd{\textcite{ariza2017labour} }\DIFaddend mostraron que la informalidad aumenta la desigualdad del ingreso laboral en las ciudades. Sin embargo, estos trabajos calculan un único estimador sin justificar la elección, estudian un factor a la vez, usan datos anteriores a 2016 y no reportan la precisión de sus resultados. Este proyecto busca llenar ese vacío\DIFdelbegin \DIFdel{al validar primero }\DIFdelend \DIFaddbegin \DIFadd{: primero valida }\DIFaddend el estimador y después \DIFdelbegin \DIFdel{sí aplicarlo}\DIFdelend \DIFaddbegin \DIFadd{lo aplica}\DIFaddend .

\section{Pregunta problema}

\begin{quote}
\emph{De cuatro estimadores candidatos del índice de Theil ---ponderado (A), sin ponderar (B), lognormal por momentos (C) y lognormal por máxima verosimilitud (D)---, ¿cuál tiene mejores propiedades estadísticas (menor sesgo, menor ECM y consistencia) para estimar la desigualdad del ingreso laboral a partir de la GEIH?}
\end{quote}


\section{Objetivos}

\subsection{Objetivo general}

Comparar las propiedades estadísticas de cuatro estimadores del índice de Theil mediante un experimento de simulación con la GEIH, y aplicar el estimador validado para medir y explicar la desigualdad del ingreso laboral en Colombia entre 2018 y 2026.

\subsection{Objetivos específicos}

\begin{enumerate}[leftmargin=0.75in]
  \item Definir formalmente los cuatro estimadores candidatos (A, B, C y D), sus supuestos y las propiedades teóricas esperadas de cada uno (sesgo, consistencia, eficiencia y suficiencia).
  \item Construir, a partir de la GEIH 2024 completa, un valor de referencia (``valor verdadero'') del Theil contra el cual evaluar los cuatro estimadores.
  \item Diseñar y ejecutar un experimento de simulación con submuestras de distintos tamaños para estimar el sesgo, la varianza, el ECM y la eficiencia relativa de cada estimador.
  \item Determinar, con base en la evidencia del experimento, cuál de los cuatro estimadores es consistente y debe usarse en el análisis aplicado.
  \item Construir \DIFdelbegin \DIFdel{nueve }\DIFdelend bases anuales comparables (\DIFdelbegin \DIFdel{enero--julio}\DIFdelend \DIFaddbegin \DIFadd{enero--junio de cada año}\DIFaddend , 2018--2026) y calcular con el estimador validado el Theil de cada año\DIFaddbegin \DIFadd{, excepto 2020, }\DIFaddend y su descomposición entre/dentro de grupos para las cinco dimensiones.
  \item Estimar intervalos de confianza del 95~\% por bootstrap y contrastar las hipótesis H1, H2 y H3.
\end{enumerate}

\section{Estimadores que pondremos a prueba}

Sea $y_i$ el ingreso laboral de la persona $i$ y $w_i$ su factor de expansión, para $i = 1, \dots, n$. El parámetro de interés es el índice de Theil de la población de ocupados:
\begin{equation}
  T = \frac{1}{N}\sum_{j=1}^{N} \frac{Y_j}{\mu_Y}\ln\!\left(\frac{Y_j}{\mu_Y}\right),
\end{equation}
donde $N$ es el tamaño de la población, $Y_j$ el ingreso de cada persona y $\mu_Y$ su media.

\textbf{\DIFdelbegin \DIFdel{Estimador }\DIFdelend \DIFaddbegin \DIFadd{Estimadores por sustitución (}\DIFaddend A \DIFaddbegin \DIFadd{y B)}\DIFaddend .} Reemplazan la distribución poblacional por la muestral. El estimador \textbf{A} usa los factores de expansión, con $W = \sum_i w_i$ y $\hat{\mu}_w = \frac{1}{W}\sum_i w_i y_i$:
\begin{equation}
  \hat{T}_A = \frac{1}{W}\sum_{i=1}^{n} w_i \frac{y_i}{\hat{\mu}_w}\ln\!\left(\frac{y_i}{\hat{\mu}_w}\right).
\end{equation}
El estimador \textbf{B} \DIFdelbegin \DIFdel{, solo }\DIFdelend es el estimador A \DIFdelbegin \DIFdel{pero }\DIFdelend sin el factor de expansión, \DIFdelbegin \DIFdel{por lo que es }\DIFdelend \DIFaddbegin \DIFadd{es decir, }\DIFaddend la misma fórmula con $w_i = 1$ para todos. Como la GEIH no le da a todas las personas la misma probabilidad de ser encuestadas, B mide la desigualdad \emph{de los encuestados} y no la de la población; por eso se espera que sea sesgado incluso con muestras muy grandes.

\textbf{Estimadores lognormales (C y D).} Si $\ln y \sim N(m, \sigma^2)$, entonces $T = \sigma^2/2$ \DIFdelbegin \DIFdel{(Cowell, 2011)}\DIFdelend \DIFaddbegin \parencite{cowell2011measuring}\DIFaddend , así que basta con estimar $\sigma^2$:
\begin{equation}
  \hat{T}_C = \frac{1}{2}\ln\!\left(1 + \frac{s_y^2}{\bar{y}^{\,2}}\right),
  \qquad
  \hat{T}_D = \frac{1}{2}\cdot\frac{1}{n}\sum_{i=1}^{n}\left(\ln y_i - \overline{\ln y}\right)^2 .
\end{equation}
C despeja $\sigma^2$ igualando la media y la varianza muestrales del ingreso con las de una lognormal (método de los momentos). D usa el estimador de máxima verosimilitud \DIFaddbegin \DIFadd{de $\sigma^2$, que depende de los datos solo a través de $\sum \ln y_i$ y $\sum (\ln y_i)^2$}\DIFaddend .


\section{Selección y exploración de bases de datos}

\subsection{Fuente}

Se usa la \textbf{Gran Encuesta Integrada de Hogares (GEIH)} del \DIFdelbegin \DIFdel{DANE (2026)}\DIFdelend \DIFaddbegin \DIFadd{\textcite{dane2026geih}}\DIFaddend : encuesta mensual, probabilística y de cobertura nacional, con una muestra de cerca de 315.000 hogares al año \DIFdelbegin \DIFdel{(DANE, 2023)}\DIFdelend \DIFaddbegin \parencite{dane2023metodologia}\DIFaddend . Es la única fuente que reúne ingreso laboral y factores de expansión\DIFdelbegin \DIFdel{, esto es realmente útil }\DIFdelend \DIFaddbegin \DIFadd{. Esto es clave, }\DIFaddend porque sin factores de expansión no se podría construir el estimador A ni evaluar el sesgo de B. Para cada año se unen los meses de enero a \DIFdelbegin \DIFdel{julio}\DIFdelend \DIFaddbegin \DIFadd{junio de tres módulos (ocupados, características generales y datos del hogar y la vivienda)}\DIFaddend , lo que da \DIFdelbegin \DIFdel{9 bases comparables}\DIFdelend \DIFaddbegin \DIFadd{nueve bases anuales (2018--2026)}\DIFaddend . Cada una tiene \textbf{más de \DIFdelbegin \DIFdel{50.000 observaciones }\DIFdelend \DIFaddbegin \DIFadd{140.000 ocupados }\DIFaddend y 12 variables} (Tabla 1), por lo que cumple los requisitos.

%DIF <  ---------------- Tabla 2 ----------------
\DIFdelbegin
\begin{table}[H]
  \singlespacing
  \raggedright
  \color{RED}{\bfseries Tabla 1 --- versión anterior\par}
  \vspace{0.5\baselineskip}
  {\itshape Variables seleccionadas de la GEIH\par}
  \vspace{0.5\baselineskip}
  \small
  \renewcommand{\arraystretch}{1.35}
  \begin{tabular}{@{}>{\raggedright\arraybackslash}p{1.45in}
                    >{\raggedright\arraybackslash}p{1.35in}
                    >{\raggedright\arraybackslash}p{1.45in}
                    >{\raggedright\arraybackslash}p{1.65in}@{}}
    \toprule
    Variable & Código GEIH & Tipo / escala & Uso \\
    \midrule
    Ingreso laboral mensual & INGLABO & Continua / razón & Cálculo del Theil \\
    Factor de expansión & FEX\_C18 & Continua / razón & Ponderación (estimador A) \\
    Nivel educativo & P3042 (antes P6210) & Cualitativa / ordinal & Dimensión 1 \\
    Cotización a pensión & P6920 & Cualitativa / nominal & Formalidad (dimensión 2) \\
    Posición ocupacional & P6430 & Cualitativa / nominal & Formalidad (dimensión 2) \\
    Estrato & P4030S1A1 & Cualitativa / ordinal & Dimensión 3 \\
    Departamento & DPTO & Cualitativa / nominal & Región (dimensión 4) \\
    Sexo & P3271 (antes P6020) & Cualitativa / nominal & Dimensión 5 \\
    Edad & P6040 & Discreta / razón & Filtro ($\geq$ 15 años) \\
    Zona & CLASE & Cualitativa / nominal & Descripción \\
    Horas trabajadas & P6800 & Discreta / razón & Descripción \\
    Rama de actividad & RAMA2D\_R4 & Cualitativa / nominal & Descripción \\
    \bottomrule
  \end{tabular}


  \vspace{0.3\baselineskip}
  {\itshape Nota.} Los códigos cambian entre los marcos muestrales 2005 y 2018 y se verificarán en el diccionario de cada año.
\end{table}
 %DIF >  ---------------- Tabla 1 ----------------
\DIFaddbegin \begin{table}[H]
  \singlespacing
  \raggedright
  \color{BLUE}{\bfseries Tabla 1 --- versión corregida\par}
  \vspace{0.5\baselineskip}
  {\itshape Variables seleccionadas de la GEIH\par}
  \vspace{0.5\baselineskip}
  \small
  \renewcommand{\arraystretch}{1.35}
  \begin{tabular}{@{}>{\raggedright\arraybackslash}p{1.45in}
                    >{\raggedright\arraybackslash}p{1.45in}
                    >{\raggedright\arraybackslash}p{1.35in}
                    >{\raggedright\arraybackslash}p{1.65in}@{}}
    \toprule
    Variable & Código GEIH & Tipo / escala & Uso \\
    \midrule
    Ingreso laboral mensual & INGLABO & Continua / razón & Cálculo del Theil \\
    Factor de expansión & FEX\_C\_2011 (2018--2021); FEX\_C18 (2022--2026) & Continua / razón & Ponderación (estimador A) \\
    Nivel educativo & P3042 (antes P6210) & Cualitativa / ordinal & Dimensión 1 \\
    Cotización a pensión & P6920 & Cualitativa / nominal & Formalidad (dimensión 2) \\
    Estrato & P4030S1A1 & Cualitativa / ordinal & Dimensión 3 \\
    Departamento & DPTO & Cualitativa / nominal & Región (dimensión 4) \\
    Sexo & P3271 (antes P6020) & Cualitativa / nominal & Dimensión 5 \\
    Posición ocupacional & P6430 & Cualitativa / nominal & Filtro (remunerados) y ajuste por no respuesta \\
    Edad & P6040 & Discreta / razón & Filtro ($\geq$ 15 años) \\
    Zona & CLASE & Cualitativa / nominal & Descripción \\
    Horas trabajadas & P6800 & Discreta / razón & Descripción \\
    Tamaño de la empresa & P3069 (antes P6870) & Cualitativa / ordinal & Descripción \\
    \bottomrule
  \end{tabular}

  \vspace{0.3\baselineskip}
  {\itshape Nota.} Los códigos cambian entre el marco muestral antiguo (2018--2021) y el Marco 2018 (2022--2026); el código guarda, para cada año, el nombre original de cada variable.
\end{table}
\DIFaddend 

\subsection{Depuración}

\DIFaddbegin \DIFadd{El principio de la depuración es }\emph{\DIFadd{marcar, no borrar}}\DIFadd{: la base de cada año conserva a todos los ocupados, y cada criterio se guarda como un indicador. Así se puede verificar cuántas personas excluye cada paso (Tabla 2). Los criterios son:
}

\DIFaddend \begin{itemize}[leftmargin=0.75in]
  \item \DIFdelbegin \DIFdel{Se conservan }\DIFdelend \DIFaddbegin \textbf{\DIFadd{Edad:}} \DIFaddend ocupados de 15 años o más\DIFdelbegin \DIFdel{con ingreso laboral mayor que cero para evitar problemas con los logaritmos}\DIFdelend \DIFaddbegin \DIFadd{.
  }\item \textbf{\DIFadd{Remunerados:}} \DIFadd{se excluyen los trabajadores sin remuneración, porque su ingreso laboral es cero por definición.
  }\item \textbf{\DIFadd{Cobertura común:}} \DIFadd{los 24 departamentos que la encuesta cubre en todos los años. El Marco 2018 agrega nueve departamentos (San Andrés y los de Orinoquía-Amazonía) que no existen en 2018--2021, así que quedan fuera para que la serie sea comparable.
  }\item \textbf{\DIFadd{Ingreso positivo:}} \DIFadd{se excluyen los ingresos iguales a cero, porque el Theil y los estimadores lognormales usan logaritmos y $\ln 0$ no está definido}\DIFaddend .
  \item \textbf{\DIFdelbegin \DIFdel{Formal}\DIFdelend \DIFaddbegin \DIFadd{Dimensiones completas}\DIFaddend :} \DIFaddbegin \DIFadd{la persona debe tener información en las cinco dimensiones.
  }\item \textbf{\DIFadd{No respuesta de ingreso:}} \DIFadd{quienes no reportan ingreso no se borran. Su peso se reparte entre quienes sí lo reportan dentro de celdas de posición ocupacional, formalidad, educación, región y sexo, de modo que el total del factor de expansión se conserva.
}\end{itemize}

\DIFadd{Las dimensiones se definen así. }\textbf{\DIFadd{Educación:}} \DIFadd{cinco niveles armonizados entre marcos (ninguno o preescolar, primaria, secundaria, media y superior). }\textbf{\DIFadd{Formalidad:}} \DIFadd{es formal quien }\DIFaddend cotiza a pensión \DIFdelbegin \DIFdel{, según su posición ocupacional. }\DIFdelend \DIFaddbegin \DIFadd{o ya está pensionado (P6920). }\textbf{\DIFadd{Estrato:}} \DIFadd{1 a 6. }\DIFaddend \textbf{Región:} \DIFdelbegin \DIFdel{regiones del DANE (}\DIFdelend \DIFaddbegin \DIFadd{cinco regiones dentro de la cobertura común (Bogotá, }\DIFaddend Caribe, Oriental, Central \DIFdelbegin \DIFdel{, }\DIFdelend \DIFaddbegin \DIFadd{y }\DIFaddend Pacífica\DIFdelbegin \DIFdel{, Bogotá, Amazonía-Orinoquía). }%DIFDELCMD < \item %%%
\DIFdel{Se usa el factor de expansión del Marco 2018 en todos los }\DIFdelend \DIFaddbegin \DIFadd{). }\textbf{\DIFadd{Sexo:}} \DIFadd{hombre o mujer.
}

\DIFadd{El año 2020 se excluye del cálculo del Theil: por la recolección telefónica durante la pandemia, falta el ingreso de la mayoría de los ocupados y no se recogió el estrato (Tabla 2). Por eso el análisis aplicado usa ocho }\DIFaddend años\DIFdelbegin \DIFdel{para que la comparación 2018--2026 sea válida.
}%DIFDELCMD < \end{itemize}
%DIFDELCMD < %%%
\DIFdelend \DIFaddbegin \DIFadd{.
}

%DIF >  ---------------- Tabla 2 ----------------
\begin{table}[H]
  \singlespacing
  \raggedright
  \color{BLUE}{\bfseries Tabla 2 --- nueva\par}
  \vspace{0.5\baselineskip}
  {\itshape Trazabilidad de la depuración por año (enero--junio)\par}
  \vspace{0.5\baselineskip}
  \small
  \renewcommand{\arraystretch}{1.25}
  \begin{tabular}{@{}lrrrr@{}}
    \toprule
    Año & Ocupados & Población de estudio & Entran al Theil & \% del FEX total \\
    \midrule
    2018 & 174.668 & 162.300 & 141.025 & 82,6 \\
    2019 & 168.441 & 156.616 & 151.283 & 87,6 \\
    2020 & 146.862 &  96.931 &  57.042 & 46,4 \\
    2021 & 146.094 & 136.479 & 129.750 & 86,2 \\
    2022 & 194.647 & 168.270 & 160.946 & 85,3 \\
    2023 & 183.136 & 160.299 & 157.776 & 88,3 \\
    2024 & 178.204 & 156.063 & 152.213 & 86,9 \\
    2025 & 176.176 & 154.659 & 153.340 & 89,1 \\
    2026 & 177.195 & 155.262 & 154.657 & 89,0 \\
    \bottomrule
  \end{tabular}

  \vspace{0.3\baselineskip}
  {\itshape Nota.} \emph{Población de estudio}: cumple todos los criterios de la depuración. \emph{Entran al Theil}: además reportan ingreso positivo. \emph{\% del FEX total}: porcentaje del factor de expansión de todos los ocupados que representan quienes entran al Theil, antes del ajuste por no respuesta. 2020 no se usa en el Theil.
\end{table}
\DIFaddend 


\section{Hipótesis}

Todas las pruebas \DIFdelbegin \DIFdel{las haremos }\DIFdelend \DIFaddbegin \DIFadd{son unilaterales }\DIFaddend con $\alpha = 0,05$.

\textbf{H1. El estimador ponderado (A) es consistente y tiene el menor ECM de los cuatro cuando el tamaño de muestra crece, mientras que B, C y D no convergen al valor verdadero.}

Para cada competidor $j$ y $n = 50.000$ se define, en cada réplica, la diferencia de pérdidas $d_j^{(r)} = (\hat{T}_A^{(r)} - T)^2 - (\hat{T}_j^{(r)} - T)^2$, y se contrasta:
\begin{equation}
  H_0:\; \bar{d}_j \geq 0 \qquad\qquad H_1:\; \bar{d}_j < 0, \qquad j \in \{B, C, D\}
\end{equation}
con una prueba $t$ \DIFaddbegin \DIFadd{pareada }\DIFaddend sobre las $R$ réplicas. Se concluye a favor de A si se rechaza $H_0$ para los tres competidores. La consistencia se evalúa además con la tendencia del ECM de cada estimador a lo largo de los cuatro tamaños de muestra.

\textbf{H2 (aplicada\DIFdelbegin \DIFdel{- esto lo hacemos una vez ya hayamos elegido }\DIFdelend \DIFaddbegin \DIFadd{, con }\DIFaddend el \DIFdelbegin \DIFdel{mejor }\DIFdelend estimador \DIFaddbegin \DIFadd{elegido}\DIFaddend ). En 2026 la educación explica más desigualdad que cada una de las otras cuatro dimensiones.}
\begin{equation}
  H_0:\; S_{edu} \leq S_k \text{ para algún } k \qquad\qquad
  H_1:\; S_{edu} > S_k \text{ para todo } k
\end{equation}

Se rechaza $H_0$ si el percentil 5 de cada diferencia bootstrap $S_{edu}^{*b} - S_k^{*b}$ es mayor que cero.

\textbf{H3 (aplicada). La desigualdad del ingreso laboral fue menor en 2026 que en 2018.}
\begin{equation}
  H_0:\; T_{2026} \geq T_{2018} \qquad\qquad
  H_1:\; T_{2026} < T_{2018}
\end{equation}

Se rechaza $H_0$ si el percentil 95 de las diferencias $T_{2026}^{*b} - T_{2018}^{*b}$ es menor que cero. Además de la decisión se reportará el tamaño de la diferencia, pues con muestras tan grandes casi todo resulta significativo.

\textbf{Limitaciones.} El valor verdadero del experimento es la propia GEIH 2024, no la población colombiana, así que la simulación evalúa los estimadores bajo un diseño que imita el real pero lo simplifica. \DIFaddbegin \DIFadd{Los factores de expansión de 2018--2021 provienen del marco muestral antiguo y los de 2022--2026 del Marco 2018, y el cuestionario también cambió; por eso parte de la diferencia entre 2018 y 2026 puede deberse al cambio metodológico y no a un cambio real de la desigualdad. El bootstrap remuestrea personas e ignora el muestreo por conglomerados, por lo que los intervalos pueden quedar algo estrechos. Las dimensiones están correlacionadas entre sí, los ingresos altos suelen subreportarse, y los resultados aplicados indican asociación, no causalidad.
}\DIFaddend 

\section{Resumen de artículos relacionados}

\subsection{\DIFdelbegin \DIFdel{Mills y Zandvakili (1997)}\DIFdelend \DIFaddbegin \texorpdfstring{\textcite{mills1997statistical}}{Mills y Zandvakili (1997)}\DIFaddend }

Proponen el bootstrap para construir intervalos y pruebas sobre el Gini y el Theil y sus descomposiciones, con datos de panel de Estados Unidos. Muestran que diferencias aparentes entre grupos pueden no ser significativas incluso con muestras grandes. Tratan el índice de desigualdad como lo que es, un estimador con distribución muestral.

\subsection{\DIFdelbegin \DIFdel{Sánchez-Torres (2017)}\DIFdelend \DIFaddbegin \texorpdfstring{\textcite{sanchez2017desigualdad}}{Sánchez-Torres (2017)}\DIFaddend }

Con la ECH y la GEIH (2002--2015), el autor descompone el Theil entre y dentro de departamentos y el Gini por fuentes de ingreso. Encuentra que cerca del 88~\% de la desigualdad ocurre dentro de los departamentos, que Chocó y La Guajira están entre los más desiguales y que el ingreso de los asalariados es la fuente que más aporta.


\subsection{\DIFdelbegin \DIFdel{Ariza y Montes-Rojas (2017)}\DIFdelend \DIFaddbegin \texorpdfstring{\textcite{ariza2017labour}}{Ariza y Montes-Rojas (2017)}\DIFaddend }

Con un panel de 13 ciudades (2002--2015), modelos de efectos fijos y variables instrumentales, los autores muestran que la informalidad aumenta significativamente la desigualdad del ingreso laboral, y que la desigualdad ocurre casi toda dentro de las ciudades.


% ---------------- Referencias ----------------
\clearpage
\section{\DIFdelbegin \DIFdel{Bibliografía}\DIFdelend \DIFaddbegin \DIFadd{Referencias}\DIFaddend }

\DIFdelbegin %DIFDELCMD < \begin{referencias}
%DIFDELCMD < \item %%%
\DIFdel{Ariza, J., y Montes-Rojas, G. (2017). Labour income inequality and the informal sector in Colombian cities. }%DIFDELCMD < \emph{Cuadernos de Economía, 36}(72)%%%
\DIFdel{, 77--98. }%DIFDELCMD < \url{https://doi.org/10.15446/cuad.econ.v36n72.65820}
%DIFDELCMD < 

%DIFDELCMD < \item %%%
\DIFdel{Cowell, F. A. (2011). }\emph{\DIFdel{Measuring inequality}} %DIFAUXCMD
\DIFdel{(3.ª ed.). Oxford University Press.
}%DIFDELCMD < 

%DIFDELCMD < \item %%%
\DIFdel{Cowell, F. A., y Jenkins, S. P. (1995). How much inequality can we explain? A methodology and an application to the United States. }%DIFDELCMD < \emph{The Economic Journal, 105}(429)%%%
\DIFdel{, 421--430. }%DIFDELCMD < \url{https://doi.org/10.2307/2235501}
%DIFDELCMD < 

%DIFDELCMD < \item %%%
\DIFdel{Departamento Administrativo Nacional de Estadística }%DIFDELCMD < [%%%
\DIFdel{DANE}%DIFDELCMD < ]%%%
\DIFdel{. (2023). }\emph{\DIFdel{Metodología general Gran Encuesta Integrada de Hogares -- GEIH}} %DIFAUXCMD
\DIFdel{(DSO-GEIH-MET-001, versión 11). }%DIFDELCMD < \url{https://www.dane.gov.co/files/operaciones/GEIH/met-GEIH.pdf}
%DIFDELCMD < 

%DIFDELCMD < \item %%%
\DIFdel{Departamento Administrativo Nacional de Estadística }%DIFDELCMD < [%%%
\DIFdel{DANE}%DIFDELCMD < ]%%%
\DIFdel{. (2025). }\emph{\DIFdel{Pobreza monetaria en Colombia: Año 2024}} %DIFAUXCMD
%DIFDELCMD < [%%%
\DIFdel{Boletín técnico}%DIFDELCMD < ]%%%
\DIFdel{. }%DIFDELCMD < \url{https://www.dane.gov.co/files/operaciones/PM/bol-PM-2024.pdf}
%DIFDELCMD < 

%DIFDELCMD < \item %%%
\DIFdel{Departamento Administrativo Nacional de Estadística }%DIFDELCMD < [%%%
\DIFdel{DANE}%DIFDELCMD < ]%%%
\DIFdel{. (2026). }\emph{\DIFdel{Gran Encuesta Integrada de Hogares -- GEIH, 2018--2026}} %DIFAUXCMD
%DIFDELCMD < [%%%
\DIFdel{Conjunto de datos}%DIFDELCMD < ]%%%
\DIFdel{. }%DIFDELCMD < \url{https://microdatos.dane.gov.co}
%DIFDELCMD < 

%DIFDELCMD < \item %%%
\DIFdel{Efron, B., y Tibshirani, R. J. (1993). }\emph{\DIFdel{An introduction to the bootstrap}}%DIFAUXCMD
\DIFdel{. Chapman \& Hall.
}%DIFDELCMD < 

%DIFDELCMD < \item %%%
\DIFdel{Mills, J. A., y Zandvakili, S. (1997). Statistical inference via bootstrapping for measures of inequality. }%DIFDELCMD < \emph{Journal of Applied Econometrics, 12}(2)%%%
\DIFdel{, 133--150.
}%DIFDELCMD < 

%DIFDELCMD < \item %%%
\DIFdel{Sánchez-Torres, R. M. (2017). Desigualdad del ingreso en Colombia: un estudio por departamentos. }%DIFDELCMD < \emph{Cuadernos de Economía, 36}(72)%%%
\DIFdel{, 139--178. }%DIFDELCMD < \url{https://doi.org/10.15446/cuad.econ.v36n72.65880}
%DIFDELCMD < 

%DIFDELCMD < \item %%%
\DIFdel{Shorrocks, A. F. (1980). The class of additively decomposable inequality measures. }%DIFDELCMD < \emph{Econometrica, 48}(3)%%%
\DIFdel{, 613--625. }%DIFDELCMD < \url{https://doi.org/10.2307/1913126}
%DIFDELCMD < 

%DIFDELCMD < \item %%%
\DIFdel{Theil, H. (1967). }\emph{\DIFdel{Economics and information theory}}%DIFAUXCMD
\DIFdel{. North-Holland.
}%DIFDELCMD < \end{referencias}
%DIFDELCMD < %%%
\DIFdelend \DIFaddbegin \printbibliography[heading=none]
\DIFaddend 

\end{document}
